\begin{examplebox}
	\textbf{\textcolor{CExample}{例1（基础）}}
	\textbf{\textcolor{CExample}{题目：}}
	已知向量 $\mathbf{a},\mathbf{b}$ 满足 $|\mathbf{a}|=3$，$|\mathbf{b}|=5$，求 $|\mathbf{a}-\mathbf{b}|$ 的取值范围，并指出取最大值和最小值时向量方向关系。
	\textbf{\textcolor{CExample}{解题步骤：}}
	\begin{description}[style=nextline, leftmargin=2.8em, labelsep=0.5em, itemsep=0.45em]
		\item[\textbf{第一步（识别不等式）}]
			由向量三角不等式：
			\[
				\bigl| |\mathbf{a}|-|\mathbf{b}| \bigr| \le |\mathbf{a}-\mathbf{b}| \le |\mathbf{a}|+|\mathbf{b}|.
			\]
			\par\textbf{理由：} 直接套用结论，注意 $\mathbf{a}-\mathbf{b}$ 对应结论中的差。
		\item[\textbf{第二步（代入数值）}]
			代入 $|\mathbf{a}|=3$，$|\mathbf{b}|=5$ 得
			\[
				|3-5| = 2 \le |\mathbf{a}-\mathbf{b}| \le 3+5 = 8.
			\]
			\par\textbf{理由：} 绝对值差与和运算。
		\item[\textbf{第三步（取等条件分析）}]
			最大值 $8$ 在 $|\mathbf{a}-\mathbf{b}|=|\mathbf{a}|+|\mathbf{b}|$ 时发生，此时 $\mathbf{a}$ 与 $\mathbf{b}$ 反向共线；最小值 $2$ 在 $|\mathbf{a}-\mathbf{b}|=||\mathbf{a}|-|\mathbf{b}||$ 时发生，此时 $\mathbf{a}$ 与 $\mathbf{b}$ 同向共线。
			\par\textbf{理由：} 差形式的三角不等式，右边取等要求 $\mathbf{a}$ 与 $-\mathbf{b}$ 同向（即 $\mathbf{a}$ 与 $\mathbf{b}$ 反向），左边取等要求 $\mathbf{a}$ 与 $-\mathbf{b}$ 反向（即 $\mathbf{a}$ 与 $\mathbf{b}$ 同向）。
	\end{description}
	\textbf{\textcolor{CExample}{关键结论：}}
	$|\mathbf{a}-\mathbf{b}|\in[2,8]$，最大值对应 $\mathbf{a},\mathbf{b}$ 反向，最小值对应同向。

	\medskip
	\textbf{\textcolor{CExample}{例2（稍变形）}}
	\textbf{\textcolor{CExample}{题目：}}
	已知向量 $\mathbf{a},\mathbf{b}$ 满足 $|\mathbf{a}|=4$，$|\mathbf{b}|=6$，且 $|\mathbf{a}+\mathbf{b}|=10$，求 $|\mathbf{a}-\mathbf{b}|$ 的值。
	\textbf{\textcolor{CExample}{解题步骤：}}
	\begin{description}[style=nextline, leftmargin=2.8em, labelsep=0.5em, itemsep=0.45em]
		\item[\textbf{第一步（识别取等）}]
			由三角不等式 $|\mathbf{a}+\mathbf{b}|\le|\mathbf{a}|+|\mathbf{b}|$，已知右端 $10$ 等于左端最大值，故取等号。
			\par\textbf{理由：} 三角不等式右边等号成立条件为 $\mathbf{a},\mathbf{b}$ 同向共线。
		\item[\textbf{第二步（同向推出差模）}]
			当 $\mathbf{a},\mathbf{b}$ 同向时，$|\mathbf{a}-\mathbf{b}| = \bigl| |\mathbf{a}|-|\mathbf{b}| \bigr| = |4-6| = 2$。
			\par\textbf{理由：} 同向向量差的模等于模之差的绝对值。
		\item[\textbf{第三步（结果验证）}]
			检查 $|\mathbf{a}-\mathbf{b}|=2$ 是否满足三角不等式：$||\mathbf{a}|-|\mathbf{b}||\le2\le|\mathbf{a}|+|\mathbf{b}|$，成立。
			\par\textbf{理由：} 一致性自检。
	\end{description}
	\textbf{\textcolor{CExample}{关键结论：}}
	$|\mathbf{a}-\mathbf{b}|=2$。

	\medskip
	\textbf{\textcolor{CExample}{例3（综合应用）}}
	\textbf{\textcolor{CExample}{题目：}}
	已知向量 $\mathbf{a},\mathbf{b},\mathbf{c}$ 满足 $|\mathbf{a}|=|\mathbf{b}|=|\mathbf{c}|=1$，求 $|\mathbf{a}+\mathbf{b}+\mathbf{c}|$ 的最大值，并说明取得最大值时向量方向关系。
	\textbf{\textcolor{CExample}{解题步骤：}}
	\begin{description}[style=nextline, leftmargin=2.8em, labelsep=0.5em, itemsep=0.45em]
		\item[\textbf{第一步（两向量求和）}]
			先对 $\mathbf{a},\mathbf{b}$ 应用三角不等式：$|\mathbf{a}+\mathbf{b}| \le |\mathbf{a}|+|\mathbf{b}| = 2$，等号当 $\mathbf{a},\mathbf{b}$ 同向时成立。
			\par\textbf{理由：} 三角不等式右边形式。
		\item[\textbf{第二步（再与第三向量和）}]
			将 $\mathbf{a}+\mathbf{b}$ 看作一个向量，对 $\mathbf{a}+\mathbf{b}$ 与 $\mathbf{c}$ 再次应用三角不等式：
			\[
				|(\mathbf{a}+\mathbf{b})+\mathbf{c}| \le |\mathbf{a}+\mathbf{b}| + |\mathbf{c}|.
			\]
			结合第一步得 $|\mathbf{a}+\mathbf{b}+\mathbf{c}| \le 2 + 1 = 3$，等号当 $|\mathbf{a}+\mathbf{b}|=2$ 且 $\mathbf{a}+\mathbf{b}$ 与 $\mathbf{c}$ 同向时成立。
			\par\textbf{理由：} 三角不等式叠加。
		\item[\textbf{第三步（等号传递整合）}]
			等号同时成立要求 $\mathbf{a},\mathbf{b}$ 同向且 $\mathbf{a}+\mathbf{b}$ 与 $\mathbf{c}$ 同向，即 $\mathbf{a},\mathbf{b},\mathbf{c}$ 三者同向共线。故最大值为 $3$。
			\par\textbf{理由：} 等号条件传递性。
	\end{description}
	\textbf{\textcolor{CExample}{关键结论：}}
	最大值 $3$，当 $\mathbf{a},\mathbf{b},\mathbf{c}$ 同向共线时取得。
\end{examplebox}
